\frametitle{Three Objectives, and Which One Is Used}
\footnotesize

A classical smoother is fixed by a few discrete choices --- a scheme, a damping factor, a
sweep count --- and analysed through the propagator $S=I-BA$. A learned smoother is fixed
instead by the weights of a network, and a network is not a scheme. The decision that
replaces the choice of scheme is therefore the choice of \emph{objective}, and it is not
unique.

\vspace{1.5mm}
\begin{tabular}{@{}p{3.5cm}p{7.1cm}p{13.2cm}@{}}
\toprule
Objective & What is minimised & What it buys, and what it costs \\
\midrule
\textbf{1. Correction matching} &
$\frac1N\sum_j\bigl\|\mathcal{C}(\kappa,r^{(j)})-Br^{(j)}\bigr\|^2$ &
Regression onto an operation that is already implemented, so the easiest to train --- and
it \emph{cannot exceed the smoother it imitates}. \\[3pt]
\textbf{2. Error damping} &
$\frac1N\sum_j\bigl\|e^{(j)}-\mathcal{C}\bigl(\kappa,Ae^{(j)}\bigr)\bigr\|^2
 \big/ \bigl\|e^{(j)}\bigr\|^2$ &
Trains the network to \emph{reduce the error} rather than to imitate an operator.
\alert{This is the objective used here.} \\[3pt]
\textbf{3. Cycle unrolling} &
$\bigl\|\mathcal{V}_\theta(0)-A^{-1}b\bigr\|$ &
Optimises the quantity ultimately of interest through a differentiable V-cycle. Most
aligned with the goal, \emph{hardest to interpret}: a gain cannot be assigned to any one
node. \\
\bottomrule
\end{tabular}

\vspace{2mm}
Objective~3 is what a complete solver would be trained against, and Stage~III will have to
answer for it \cite{huang2021learning}. Objective~2 is what Stage~I can \emph{interpret}:
it measures \emph{one step}, so a result is attributable to that step and not to a
composition.

\bre
The choice is a modelling decision, not a technicality. Objective~2 is also the only one of
the three that can be scored on the same test errors, in the same norm, for a network and
for Jacobi, Gauss--Seidel and SSOR alike --- which is what makes the classical baseline
fair.
\ere
